\documentclass[11pt]{article}

\usepackage[margin=1in]{geometry}
\usepackage[T1]{fontenc}
\usepackage{lmodern}
\usepackage{amsmath,amssymb,amsthm}
\usepackage{bm}
\usepackage{booktabs}
\usepackage{tabularx}
\usepackage{array}
\usepackage{graphicx}
\usepackage{siunitx}
\usepackage{xcolor}
\usepackage{enumitem}
\usepackage{microtype}
\usepackage[round,numbers,sort&compress]{natbib}
\usepackage[colorlinks=true,linkcolor=blue,citecolor=blue,urlcolor=blue]{hyperref}

% ---------------------------------------------------------------- shortcuts
\newcommand{\R}{\mathbb{R}}
\newcommand{\pdd}[2]{\frac{\partial #1}{\partial #2}}
\newcommand{\pddt}[2]{\frac{\partial^{2} #1}{\partial #2^{2}}}
\newcommand{\erfc}{\operatorname{erfc}}
\newcommand{\Pe}{\mathrm{Pe}}
\newcommand{\LL}{\mathcal{L}}
\newcommand{\RR}{\mathcal{R}}
\newcommand{\DD}{\mathbf{D}}
\newcommand{\vv}{\mathbf{v}}
\newcommand{\xx}{\mathbf{x}}
\newcommand{\nn}{\mathbf{n}}
\newcommand{\norm}[1]{\left\lVert #1 \right\rVert}
\newcommand{\abs}[1]{\left\lvert #1 \right\rvert}
\newcommand{\intO}{\int_{\Omega}}
\newcommand{\intG}{\int_{\Gamma_{\mathrm{out}}}}
\newcommand{\dd}{\,\mathrm{d}}
\newcommand{\T}{\mathsf{T}}
\newcommand{\tabhead}[1]{\textbf{#1}}
\newcolumntype{Y}{>{\raggedright\arraybackslash}X}

\title{From One to Three Dimensions\\[2pt]
\large A staged plan for groundwater contaminant transport\\
with PINNs and Fourier neural operators}
\author{Sergio Lucero}
\date{\today}

\begin{document}
\maketitle

\begin{abstract}
Phase~1 of the debugging plan established a trusted one-dimensional
advection--dispersion sandbox: the physics-informed neural network (PINN)
matches the Ogata--Banks solution to a relative $L^{2}$ error of
$1.35\times10^{-2}$ on the full grid and $8.0\times10^{-3}$ in the interior,
with the remaining error confined to the measure-zero inlet corner.  This note
explains the companion documents \texttt{2D\_plan.md} and
\texttt{3D\_plan.md}: how the governing equations, reference solutions,
network architectures, operator-learning strategies, metrics, and compute
budgets change as the model is extended to two and three spatial dimensions.
The central principle is gate-driven development: each stage adds one physical
ingredient, is checked against a closed-form, manufactured, or
grid-converged reference, and must pass an explicit accuracy threshold before
the next stage begins.  Two dimensions are required for transverse dispersion,
well capture, and heterogeneous conductivity; three dimensions are justified
only by vertical flow, partially penetrating wells, vertical heterogeneity, or
density effects, and are approached through a 2.5D axisymmetric model first.
\end{abstract}

\tableofcontents

\section{Context and scope}
\label{sec:context}

The project evolves in four logical layers:
\begin{enumerate}[leftmargin=1.6em,itemsep=2pt]
  \item a \emph{physical model} of flow and contaminant transport;
  \item \emph{reference solutions} against which every learned model is judged;
  \item \emph{learned solvers}: PINNs for a single scenario and Fourier neural
        operators (FNOs) for a family of scenarios;
  \item a \emph{verification and acceptance protocol} that prevents a
        plausible-looking but wrong result from advancing.
\end{enumerate}
The one-dimensional Phase~1 sandbox already exercised all four layers in the
simplest setting.  The 2D and 3D plans keep the same layers and add one new
physical ingredient at a time.

\paragraph{Lessons carried forward from Phase~1.}
\begin{itemize}[leftmargin=1.6em,itemsep=2pt]
  \item \textbf{Non-dimensionalize first.}  The dimensionless ADE keeps the
        residual and every loss term $\mathcal{O}(1)$.
  \item \textbf{Report the interior metric.}  The step-inlet problem is
        discontinuous at $(\xi,\tau)=(0,0)$; the full-grid error is dominated
        by a measure-zero corner.  Two and three dimensions introduce
        additional inlet/wall and well singularities, so an interior or
        singularity-excluded metric must always accompany the full-grid one.
  \item \textbf{Verify the graph before the physics.}  An extruded 1D test is
        the cheapest way to prove that a 2D or 3D implementation has not
        broken the trusted result.
  \item \textbf{Decouple flow and transport first.}  Freezing the velocity
        field removes the second derivatives of the head from the transport
        residual and makes loss balancing tractable.  A monolithic
        flow--transport PINN is a later experiment.
  \item \textbf{One gate at a time.}  A stage is ``done'' only when its
        reference comparison passes a stated threshold; see
        Sections~\ref{sec:gates}--\ref{sec:3d}.
\end{itemize}

\paragraph{Relationship of the two plans.}
The 2D plan is the primary path.  The 3D plan is conditional: it is entered
only after the corresponding 2D gates pass, and its first stage is a 2.5D
axisymmetric or vertical-slice model.  If a 2.5D model answers the site
question, the full 3D model is unnecessary; the justification for using 3D
must be written down before the first full 3D run.

\section{Mathematical model}
\label{sec:model}

\subsection{Flow}

For a confined, constant-density aquifer, the steady flow problem is
\begin{equation}
\label{eq:flow}
-\nabla\cdot(K\nabla h)=W,
\qquad
\mathbf{q}=-K\nabla h,
\qquad
\vv=\frac{\mathbf{q}}{\phi},
\end{equation}
where $h$ is hydraulic head, $K$ hydraulic conductivity, $\mathbf{q}$ specific
discharge, $\phi$ porosity, $\vv$ seepage (pore) velocity, and $W$ a
source/sink term (extraction negative, injection positive).  In 2D one usually
works with the transmissivity $T=Kb$, where $b$ is the aquifer thickness; in 3D
$K$ may be the anisotropic tensor
\begin{equation}
\label{eq:Kaniso}
K=\operatorname{diag}(K_h,K_h,K_v),
\end{equation}
with horizontal isotropy and vertical anisotropy ratio $\kappa=K_v/K_h$.

On a Dirichlet boundary $\Gamma_D$, $h=h_D$; on a no-flow boundary
$\Gamma_N$, $\mathbf{q}\cdot\nn=0$.  A well screened over an interval is
represented by distributing its rate over the screened cells rather than by a
point source; this avoids the logarithmic well singularity in the discrete
problem.

\subsection{Transport}

The conservative advection--dispersion equation (ADE) with linear equilibrium
sorption and first-order decay is
\begin{equation}
\label{eq:ade}
R\pdd{C}{t}+\nabla\cdot(\vv C)-\nabla\cdot(\DD\nabla C)+\lambda R C=S,
\end{equation}
with retardation factor $R=1+\rho_bK_d/\phi$, decay rate $\lambda\ge0$, and
source/sink $S$.  For constant $\phi$ and a divergence-free velocity field,
$\nabla\cdot(\vv C)=\vv\cdot\nabla C$.  In two dimensions the dispersion
tensor is
\begin{equation}
\label{eq:D2}
\DD=\alpha_T\abs{\vv}\mathbf{I}
+(\alpha_L-\alpha_T)\frac{\vv\vv^{\T}}{\abs{\vv}}
+D_m\tau\mathbf{I},
\end{equation}
with longitudinal dispersivity $\alpha_L$, transverse dispersivity $\alpha_T$,
molecular diffusion $D_m$, and tortuosity $\tau$.  In components, with
$\vv=(v_x,v_y)$,
\begin{equation}
\label{eq:D2comp}
\begin{aligned}
D_{xx}&=\alpha_L\abs{\vv}\frac{v_x^2}{\abs{\vv}^2}
       +\alpha_T\abs{\vv}\frac{v_y^2}{\abs{\vv}^2}+D_m\tau,\\
D_{yy}&=\alpha_L\abs{\vv}\frac{v_y^2}{\abs{\vv}^2}
       +\alpha_T\abs{\vv}\frac{v_x^2}{\abs{\vv}^2}+D_m\tau,\\
D_{xy}&=(\alpha_L-\alpha_T)\abs{\vv}\frac{v_xv_y}{\abs{\vv}^2}.
\end{aligned}
\end{equation}
In three dimensions the same construction uses a longitudinal dispersivity
$\alpha_L$, a horizontal transverse dispersivity $\alpha_{TH}$, and a vertical
transverse dispersivity $\alpha_{TV}$; the 2D form \eqref{eq:D2} is the
special case $\alpha_{TH}=\alpha_{TV}=\alpha_T$ \citep{bear1972}.

\subsection{Boundary and initial conditions}

The 1D sandbox used a step inlet and a clean initial aquifer:
\begin{equation}
\label{eq:icbc1d}
C(x,0)=0,\qquad C(0,t)=C_0,\qquad C(L,t)=0.
\end{equation}
The 2D and 3D tests keep the same structure but add lateral and vertical
boundaries:
\begin{itemize}[leftmargin=1.6em,itemsep=2pt]
  \item inlet $\xi=0$: Dirichlet $C=C_0$ (or a specified source function);
  \item outlet $\xi=1$: Dirichlet $C=0$ for the semi-infinite verification
        problems, or a zero-diffusive-flux/advective outflow condition for
        finite domains;
  \item lateral and bottom/top boundaries: no-flux,
        $\DD\nabla C\cdot\nn=0$ (or $\vv\cdot\nn=0$ when advection dominates);
  \item initial condition: $C(\xx,0)=C_0(\xx)$, usually zero or a source-zone
        blob.
\end{itemize}
For a point source, the source term is
$S=\dot M\,\delta(\xx-\xx_s)$ in \eqref{eq:ade}, with $\dot M$ the mass release
rate; for a finite source zone, $S$ is distributed over the zone.

\subsection{Non-dimensionalization}

Introduce
\begin{equation}
\label{eq:nondim}
\xi=\frac{x}{L},\quad \eta=\frac{y}{L},\quad \zeta=\frac{z}{H},\quad
\tau=\frac{tU}{L},\quad u=\frac{C}{C_0},\quad
\tilde h=\frac{h}{H_0},\quad \tilde K=\frac{K}{K_0},
\end{equation}
and the dimensionless groups
\begin{equation}
\label{eq:groups}
\Pe_L=\frac{UL}{D_L},\qquad
\Pe_T=\frac{UL}{D_T},\qquad
\Pe_{TH}=\frac{UL}{D_{TH}},\qquad
\Pe_{TV}=\frac{UH}{D_{TV}},\qquad
\kappa=\frac{K_v}{K_h},\qquad
\varepsilon=\frac{H}{L}.
\end{equation}
The transport equation becomes
\begin{equation}
\label{eq:adend}
R\pdd{u}{\tau}+\tilde{\vv}\cdot\nabla_{\xi}u
-\frac{1}{\Pe_L}\nabla_{\xi}\cdot(\tilde{\DD}\nabla_{\xi}u)
+\tilde\lambda R u=\tilde S,
\end{equation}
with $\tilde{\DD}=\DD/D_L$, $\tilde\lambda=\lambda L/U$, and
$\tilde S=SL/(C_0U)$.  The flow equation becomes
\begin{equation}
\label{eq:flownd}
-\nabla_{\xi}\cdot(\tilde K\nabla_{\xi}\tilde h)=\tilde W,
\qquad
\tilde W=\frac{WL^2}{K_0H_0}.
\end{equation}
Every quantity in \eqref{eq:adend}--\eqref{eq:flownd} is $\mathcal{O}(1)$,
which is essential for PINN training at high P\'eclet numbers.

\section{Reference solutions and the verification hierarchy}
\label{sec:references}

No learned model should be trusted without a reference.  The 2D/3D plans use
four kinds of references, in increasing order of physical realism.

\subsection{Extruded one-dimensional solution}

For uniform flow $\vv=(U,0)$ (or $(U,0,0)$), an inlet and initial condition
independent of the transverse coordinates, and no-flux lateral boundaries, the
exact solution is the Ogata--Banks solution extruded in the transverse
directions:
\begin{equation}
\label{eq:extruded}
u_{\mathrm{ref}}(\xi,\eta,\zeta,\tau)=u_{\mathrm{OB}}(\xi,\tau)
=\frac12\left[
\erfc\!\left(\frac{(\xi-\tau)\sqrt{\Pe_L}}{2\sqrt{\tau}}\right)
+e^{\Pe_L\xi}\erfc\!\left(\frac{(\xi+\tau)\sqrt{\Pe_L}}{2\sqrt{\tau}}\right)
\right].
\end{equation}
This is the first test of any 2D or 3D implementation: if the code cannot
reproduce \eqref{eq:extruded}, the new dimensions, masks, or autodiff graph are
wrong.

\subsection{Green's functions for point and line sources}

For a continuous point source in an infinite uniform-flow aquifer, the steady
two-dimensional solution is
\begin{equation}
\label{eq:green2d}
C_{\mathrm{ss}}(x,y)=\frac{Q}{2\pi b\sqrt{D_LD_T}}
\exp\!\left(\frac{Ux}{2D_L}\right)
K_0\!\left(\frac{U}{2D_L}\sqrt{x^2+\frac{D_L}{D_T}y^2}\right),
\end{equation}
where $Q$ is the mass release rate, $b$ the aquifer thickness, and $K_0$ the
modified Bessel function of the second kind.  The transient continuous-source
solution is the convolution of the instantaneous Green's function:
\begin{equation}
\label{eq:green2dt}
C(x,y,t)=\int_0^t
\frac{\dot M}{4\pi\phi b s\sqrt{D_LD_T}}
\exp\!\left[-\frac{(x-Us)^2}{4D_Ls}-\frac{y^2}{4D_Ts}\right]\dd s.
\end{equation}
In three dimensions, the steady point-source solution is
\begin{equation}
\label{eq:green3d}
C_{\mathrm{ss}}(x,y,z)=\frac{Q}{4\pi\sqrt{D_LD_TD_V}}
\exp\!\left(\frac{Ux}{2D_L}\right)\frac{e^{-mR}}{R},
\end{equation}
with
\begin{equation}
\label{eq:green3dR}
R=\sqrt{\frac{x^2}{D_L}+\frac{y^2}{D_T}+\frac{z^2}{D_V}},
\qquad
m=\frac{U}{2\sqrt{D_L}},
\end{equation}
and the transient continuous-source solution is
\begin{equation}
\label{eq:green3dt}
C=\int_0^t
\frac{\dot M}{\phi(4\pi s)^{3/2}\sqrt{D_LD_TD_V}}
\exp\!\left[-\frac{(x-Us)^2}{4D_Ls}
-\frac{y^2}{4D_Ts}-\frac{z^2}{4D_Vs}\right]\dd s.
\end{equation}
Equations \eqref{eq:green2d}--\eqref{eq:green3dt} are implemented in
\texttt{physics/references.py} and \texttt{physics/references3d.py}.  They are
used for the source/transverse-dispersion stages and for checking plume
moments and mass balance.

\subsection{Method of images for no-flux boundaries}

For a no-flux wall at $y=0$, add an image source at $-y_s$; for a wall at
$y=L_y$, add an image at $2L_y-y_s$.  For two parallel walls, the solution is
an infinite sum of images.  The same construction applies in $z$ for a
confined aquifer.  In practice, truncate the image series when the contribution
is below the target tolerance, or use a finite-volume reference solver for the
bounded domain.  Image solutions are especially useful for testing whether the
PINN's no-flux condition is actually satisfied.

\subsection{Method of manufactured solutions}

For heterogeneous conductivity, layered media, and arbitrary boundary
conditions, no simple closed form exists.  The method of manufactured solutions
(MMS) supplies an exact solution by construction.  Choose a smooth function
$u^*(\xi,\eta,\tau)$, substitute it into the PDE, and define the source
\begin{equation}
\label{eq:mms}
S^*=R\pdd{u^*}{\tau}+\vv\cdot\nabla u^*
-\nabla\cdot(\DD\nabla u^*)+\lambda R u^*.
\end{equation}
Then $u^*$ is the exact solution of \eqref{eq:ade} with source $S^*$.  Impose
$u^*$ as the initial and Dirichlet data, and require the learned model to
recover it.  The same construction verifies the flow solver: choose
$h^*(x,y,z)$ and $K(x,y,z)$, and set
\begin{equation}
\label{eq:mmsflow}
W^*=-\nabla\cdot(K\nabla h^*).
\end{equation}
A useful MMS family is
\begin{equation}
\label{eq:mmsfamily}
u^*(\xi,\eta,\zeta,\tau)=e^{-\gamma\tau}
\sin(\pi\xi)\sin(\pi\eta)\sin(\pi\zeta),
\end{equation}
which vanishes on a box and exercises all second derivatives.  A localized
Gaussian plume is a better test of a sharp front.

\subsection{Numerical reference and grid convergence}

For random conductivity fields and realistic remediation scenarios, the
reference is a mass-conservative finite-volume (FV) solver:
\begin{itemize}[leftmargin=1.6em,itemsep=2pt]
  \item cell-centered FV with two-point flux approximation for Darcy flow and
        harmonic-mean face conductivities;
  \item upwind or flux-limited advection, implicit diffusion, and operator
        splitting or an implicit--explicit time integrator;
  \item positivity-preserving limiting at high $\Pe_L$;
  \item wells distributed over screened cells;
  \item a grid-convergence study on $64^2$, $128^2$, $256^2$ in 2D and
        $32^3$, $64^3$, $128^3$ in 3D.
\end{itemize}
The FV solution is the ground truth for FNO training and for the heterogeneous
PINN tests.  It must itself pass MMS and mass-balance tests before it is used
as a reference.

\section{PINN formulation in two and three dimensions}
\label{sec:pinn}

\subsection{Residual and loss}

The network $u_\theta$ approximates the dimensionless concentration.  The PDE
residual is
\begin{equation}
\label{eq:res}
\RR_\theta(\xx,\tau)=
R\pdd{u_\theta}{\tau}+\tilde{\vv}\cdot\nabla u_\theta
-\frac{1}{\Pe_L}\nabla\cdot(\tilde{\DD}\nabla u_\theta)
+\tilde\lambda R u_\theta-\tilde S,
\end{equation}
computed by automatic differentiation.  The composite loss is
\begin{equation}
\label{eq:loss}
\LL(\theta)=\lambda_r\LL_{\mathrm{res}}
+\lambda_b\LL_{\mathrm{bc}}
+\lambda_0\LL_{\mathrm{ic}}
+\lambda_d\LL_{\mathrm{data}}
+\lambda_{\mathrm{div}}\LL_{\mathrm{div}}
+\lambda_{\mathrm{flux}}\LL_{\mathrm{flux}},
\end{equation}
with, for example,
\begin{equation}
\label{eq:lossres}
\LL_{\mathrm{res}}=\frac1{N_r}\sum_{i=1}^{N_r}
\abs{\RR_\theta(\xx_i,\tau_i)}^2,
\end{equation}
and analogous mean-squared penalties for the boundary, initial, and data
mismatches.  For coupled flow--transport, add the flow residual
$\abs{-\nabla\cdot(K\nabla h_\theta)-W}^2$ and a divergence penalty
$\abs{\nabla\cdot\vv_\theta}^2$.  All residuals are non-dimensional.

\subsection{Hard constraints on a box}

Soft penalties are easy but can be violated.  A hard ansatz has the form
\begin{equation}
\label{eq:hard}
u_\theta(\xx,\tau)=P(\xx,\tau)+B(\xx,\tau)N_\theta(\xx,\tau),
\end{equation}
where the lift $P$ satisfies the Dirichlet data and the mask $B$ vanishes on
the Dirichlet boundaries and at $\tau=0$.  For the unit square with inlet at
$\xi=0$, outlet at $\xi=1$, and no-flux walls at $\eta=0,1$, choose
\begin{equation}
\label{eq:hard2d}
s(\tau)=1-e^{-\tau/\varepsilon},\qquad
P(\xi,\eta,\tau)=(1-\xi)g(\eta,\tau)s(\tau),
\end{equation}
\begin{equation}
\label{eq:hard2dB}
B(\xi,\eta,\tau)=\xi(1-\xi)\left[1+\gamma\eta^2(1-\eta)^2\right]s(\tau).
\end{equation}
Then $u(\xi,\eta,0)=0$ exactly, $u(0,\eta,\tau)=g(\eta,\tau)s(\tau)\to g$
exponentially fast, $u(1,\eta,\tau)=0$, and
$\partial B/\partial\eta=0$ at $\eta=0,1$.  The mask respects the no-flux
walls without forcing $u=0$ there.  The three-dimensional box is the obvious
product extension:
\begin{equation}
\label{eq:hard3dB}
B=\xi(1-\xi)\left[1+\gamma\eta^2(1-\eta)^2\right]
\left[1+\gamma\zeta^2(1-\zeta)^2\right]s(\tau).
\end{equation}
For a general site geometry, replace the product mask by a signed-distance
function or an R-function; for flux boundaries, either use a mask with zero
normal derivative or enforce the flux softly and monitor the boundary residual.

\subsection{Sampling, curriculum, and optimization}

\begin{itemize}[leftmargin=1.6em,itemsep=2pt]
  \item Sample collocation points with Sobol or Latin-hypercube sequences in
        $(\xi,\eta,\tau)$ or $(\xi,\eta,\zeta,\tau)$.
  \item Oversample near the inlet, the source, the plume front, layer
        interfaces, and wells.
  \item Use residual-based adaptive refinement (RAR): every $N$ steps, add
        points where $\abs{\RR_\theta}$ is largest.
  \item For high $\Pe_L$, use a causal time weighting or train in time
        windows; a single global time interval often stalls.
  \item Optimize with Adam first (resampled collocation), then L-BFGS on a
        fixed set.  Use a P\'eclet-number curriculum: start at a moderate
        $\Pe_L$, then ramp to the target.
  \item Log per-term losses, residual percentiles, boundary/initial errors,
        gradient norms, mass balance, and wall-clock per $10^3$ steps.
\end{itemize}

\subsection{Decoupled versus monolithic flow--transport}

In the decoupled approach, the flow problem is solved once (analytically or by
FV), the velocity is evaluated on the transport collocation points, and only
the transport network is trained.  This is the recommended path for the first
2D and 3D stages.  In the monolithic approach, a shared trunk produces $h$ and
$C$, and the loss contains both residuals.  The transport residual then
contains second derivatives of $h$ through $\vv=-K\nabla h/\phi$ and through
the dispersion tensor, which makes the graph deeper and the loss balance
harder.  Attempt the monolithic version only after the decoupled version
passes its gate.

\subsection{Domain decomposition}

For layered media and large domains, split the domain into subdomains
(by layer, by vertical column, or by time window) and train one network per
subdomain with interface conditions:
\begin{equation}
\label{eq:interface}
h^{(1)}=h^{(2)},\qquad
C^{(1)}=C^{(2)},\qquad
\mathbf{q}^{(1)}\cdot\nn=\mathbf{q}^{(2)}\cdot\nn,\qquad
\DD^{(1)}\nabla C^{(1)}\cdot\nn=\DD^{(2)}\nabla C^{(2)}\cdot\nn .
\end{equation}
This is the extended PINN (XPINN) idea \citep{jagtap2020xpinn}.  It reduces the
width needed per network and localizes the difficult interfaces.  In 3D it is
also a memory strategy: each subdomain network sees a smaller batch.

\section{Fourier neural operator formulation}
\label{sec:fno}

\subsection{Operator learning}

A PINN learns one solution; an FNO learns the solution operator
\begin{equation}
\label{eq:operator}
\mathcal{G}_\theta:\;a\;\longmapsto\;u,
\end{equation}
where $a$ collects the input functions and parameters: conductivity $K$, source
location and strength, well rates and screens, and boundary data.  The output
is the concentration field $u(\xx,\tau)$, or a time sequence of fields.  The FNO
is a composition of a lifting multilayer perceptron (MLP), $L$ Fourier layers,
and a projection MLP \citep{li2021fno,kovachki2023neural}.

\subsection{Spectral convolution}

A Fourier layer applies a linear operator in frequency space:
\begin{equation}
\label{eq:spectral}
(\mathcal{K}v)(\xx)=\mathcal{F}^{-1}\!\left[
R\odot\mathcal{F}(v)\right](\xx),
\end{equation}
where $R$ is a learnable complex weight tensor truncated to $k_{\max}$ modes in
each dimension.  In 2D, $R\in\mathbb{C}^{k_x\times k_y\times c_{\mathrm{in}}
\times c_{\mathrm{out}}}$; in 3D,
$R\in\mathbb{C}^{k_x\times k_y\times k_z\times c_{\mathrm{in}}\times
c_{\mathrm{out}}}$.  The parameter count per layer is
\begin{equation}
\label{eq:params}
N_{\mathrm{spec}}=
\begin{cases}
k_xk_y\,c_{\mathrm{in}}c_{\mathrm{out}}, & \text{2D},\\[2pt]
k_xk_yk_z\,c_{\mathrm{in}}c_{\mathrm{out}}, & \text{3D}.
\end{cases}
\end{equation}
For example, $8^3$ modes with 32 channels give $5.2\times10^5$ parameters per
layer, while $16^3$ modes with 64 channels give $1.7\times10^7$ per layer.
This is why 3D FNOs start with 8--12 modes and 32 channels, and why factorized
or low-rank spectral convolutions are useful at $128^3$.

\subsection{Autoregressive rollout and loss}

For long remediation horizons, use a 2D or 3D FNO in space and roll it out in
time:
\begin{equation}
\label{eq:rollout}
C^{n+1}=\mathcal{G}_\theta(C^n,K,S;\Delta\tau),
\qquad n=0,\dots,N_t-1,
\end{equation}
with the rollout loss
\begin{equation}
\label{eq:rollloss}
\LL_{\mathrm{roll}}=\sum_{n=1}^{N_t}\omega_n
\norm{C^n-C^n_{\mathrm{ref}}}_2^2
+\lambda_{\mathrm{one}}\norm{C^1-C^1_{\mathrm{ref}}}_2^2 .
\end{equation}
The one-step term stabilizes early training; the rollout term controls error
growth.  Use pushforward or temporal bundling to avoid the optimizer exploiting
a single-step shortcut.  Always evaluate on a horizon longer than the training
horizon.

\subsection{Physics-informed FNO}

If data are scarce, add a residual penalty computed on the output grid by
finite differences or by spectral derivatives:
\begin{equation}
\label{eq:fnopde}
\LL_{\mathrm{pde}}=\sum_{n}\norm{\RR[C^n]}_2^2 .
\end{equation}
This is cheaper than autodiff through the FNO and is a useful regularizer.
It does not replace the reference data for verification.

\section{Metrics and acceptance gates}
\label{sec:gates}

\subsection{Field errors}

The primary metric is the relative $L^2$ error
\begin{equation}
\label{eq:relL2}
E_{\mathrm{rel}}=\frac{\norm{u_\theta-u_{\mathrm{ref}}}_2}
{\norm{u_{\mathrm{ref}}}_2},
\qquad
\norm{f}_2^2=\frac1{N}\sum_{i=1}^{N}f(\xx_i)^2,
\end{equation}
evaluated on a fixed grid independent of the training points.  Always report a
second, singularity-excluded metric
\begin{equation}
\label{eq:relL2int}
E_{\mathrm{rel}}^{\mathrm{int}}
=\frac{\norm{u_\theta-u_{\mathrm{ref}}}_{2,\Omega_{\mathrm{int}}}}
{\norm{u_{\mathrm{ref}}}_{2,\Omega_{\mathrm{int}}}},
\end{equation}
where $\Omega_{\mathrm{int}}$ excludes the inlet corner, inlet/wall edges, and
well neighborhoods.  In 2D and 3D the singular sets are larger than in 1D, so
this is not optional.

\subsection{Mass balance}

For a conservative model, define
\begin{equation}
\label{eq:mass}
M(t)=\intO\phi R C\dd V,
\end{equation}
and the relative balance residual
\begin{equation}
\label{eq:massbal}
\mathcal{B}(t)=
\frac{\left|M(t)-M(0)-\int_0^t\left[
\intO S\dd V-\intG\phi\vv\cdot\nn C\dd A
-\intO\lambda R\phi C\dd V\right]\dd t'\right|}
{\max\{M(t),M_{\mathrm{ref}}\}}.
\end{equation}
A model that fits the field but violates \eqref{eq:massbal} is not a valid
transport model.  Report $\max_t\mathcal{B}(t)$ and the time-averaged value.

\subsection{Breakthrough, cleanup time, and plume moments}

At a monitoring well $w$, let $C_w(t)$ be the breakthrough curve.  Report
\begin{equation}
\label{eq:btc}
E_{\mathrm{BTC}}=\frac{\norm{C_w-C_{w,\mathrm{ref}}}_2}
{\norm{C_{w,\mathrm{ref}}}_2},
\qquad
t_a=\inf\{t:C_w(t)\ge C_{\mathrm{thr}}\},
\qquad
t_c=\inf\{t:C_w(t)\le C_{\mathrm{lim}}\},
\end{equation}
where $t_a$ is the arrival time and $t_c$ the cleanup time.  Plume moments are
\begin{equation}
\label{eq:moments}
M_0=\intO\phi R C\dd V,\qquad
\xx_c=\frac1{M_0}\intO\xx\,\phi R C\dd V,
\end{equation}
\begin{equation}
\label{eq:cov}
\boldsymbol{\Sigma}=\frac1{M_0}\intO
(\xx-\xx_c)(\xx-\xx_c)^{\T}\phi R C\dd V .
\end{equation}
Compare the total mass, centroid, and covariance (spread) to the reference;
these catch errors that a pointwise $L^2$ metric can miss.

\subsection{FNO-specific metrics}

For an FNO, report:
\begin{enumerate}[leftmargin=1.6em,itemsep=2pt]
  \item one-step error and the error at $10$, $50$, and $100$ rollout steps;
  \item error versus horizon (growth rate and stability);
  \item spectral error (energy in the wrong modes);
  \item resolution-transfer error: train on a coarse grid, evaluate on a fine
        grid;
  \item conservation and positivity violations;
  \item inference cost and speedup relative to the FV reference.
\end{enumerate}

\subsection{Gate tables}

\begin{table}[htbp]
\centering
\small
\begin{tabularx}{\textwidth}{@{}lYl@{}}
\toprule
\tabhead{Gate} & \tabhead{Check} & \tabhead{Threshold}\\
\midrule
G2.1 & Extruded 1D, interior rel-$L^2$ & $\le 1.5\times10^{-2}$\\
G2.1 & Wall-normal gradient $\norm{\partial u/\partial\eta}$ & $\le 10^{-3}$\\
G2.2 & Plume-core rel-$L^2$ & $\le 5\%$\\
G2.2 & Mass-balance residual & $\le 1\%$\\
G2.3 & Head / velocity rel-$L^2$ (away from wells) & $\le 10^{-3}$ / $\le 10^{-2}$\\
G2.4 & FV observed order / PINN MMS / FNO held-out $K$ & $\ge 1.8$ / $\le 2\times10^{-2}$ / $\le 5\%$\\
G2.5 & Cleanup-time error / FNO rollout / speedup & $\le 10\%$ / $\le 5\%$ / $\ge 100\times$\\
\bottomrule
\end{tabularx}
\caption{Two-dimensional acceptance gates.  A gate is a go/no-go condition,
not a suggestion.}
\label{tab:gates2d}
\end{table}

\begin{table}[htbp]
\centering
\small
\begin{tabularx}{\textwidth}{@{}lYl@{}}
\toprule
\tabhead{Gate} & \tabhead{Check} & \tabhead{Threshold}\\
\midrule
G3.1 & Extruded 1D/2D, interior rel-$L^2$ & $\le 1.5\times10^{-2}$\\
G3.1 & $\norm{\partial u/\partial\eta}$, $\norm{\partial u/\partial\zeta}$ & $\le 10^{-3}$\\
G3.2 & 3D Green's function, plume core & $\le 5\%$\\
G3.2 & Mass-balance residual & $\le 1\%$\\
G3.3 & 3D vs axisymmetric head / velocity & $\le 2\%$ / $\le 5\%$\\
G3.4 & FV observed order / PINN MMS / FNO held-out $K$ & $\ge 1.8$ / $\le 2\times10^{-2}$ / $\le 10\%$\\
G3.5 & FNO one-step / 100-step rollout & $\le 2\%$ / $\le 10\%$\\
G3.5 & Resolution transfer / conservation & $\le 15\%$ / $\le 2\%$\\
G3.6 & Cleanup time / extracted mass / speedup & $\le 10\%$ / $\le 10\%$ / $\ge 100\times$\\
\bottomrule
\end{tabularx}
\caption{Three-dimensional acceptance gates.  G3.3 is the first gate that
requires a real 3D flow field; everything before it can be checked with
axisymmetric or extruded references.}
\label{tab:gates3d}
\end{table}

\section{Two-dimensional roadmap}
\label{sec:2d}

The 2D plan is organized so that each stage isolates one new effect.

\begin{table}[htbp]
\centering
\small
\begin{tabularx}{\textwidth}{@{}lYl@{}}
\toprule
\tabhead{Phase} & \tabhead{Work} & \tabhead{Gate}\\
\midrule
2D-0 & Refactor into \texttt{physics/}; port config, manifest, checkpoint
       helpers; add 2D MMS unit tests. & tests pass\\
2D-1 & Uniform flow, extruded 1D verification. & G2.1\\
2D-2 & Point/line source, transverse dispersion, anisotropy
       $\alpha_L/\alpha_T=1,10,100$. & G2.2\\
2D-3 & Darcy flow, extraction/injection wells; decoupled velocity first. & G2.3\\
2D-4 & Layered and random log-normal $K$; MMS; FV reference and grid
       convergence. & G2.4\\
2D-5 & Pump-and-treat/PRB scenario; FNO2d; rollout; surrogate. & G2.5\\
2D-6 & Report, figures, update \texttt{Actions.md} and \texttt{HUMAN.md}. & all gates\\
\bottomrule
\end{tabularx}
\caption{Two-dimensional work breakdown.  Do not start 2D-4 before G2.2, and do
not start 2D-5 before G2.3 and the FV solver are trusted.}
\label{tab:2d}
\end{table}

\paragraph{Why this order?}
Stage 2D-1 proves the implementation.  Stage 2D-2 introduces the first genuinely
two-dimensional physics (transverse spreading) and has a closed-form reference.
Stage 2D-3 introduces wells and a spatially varying velocity; decoupling keeps
the transport graph manageable.  Stage 2D-4 introduces heterogeneity, which
requires MMS and a numerical reference.  Stage 2D-5 is the application stage,
where the FNO surrogate becomes useful because many scenarios are needed for
uncertainty quantification and optimization.

\paragraph{Expected 2D result.}
For the extruded case, the interior error should be statistically
indistinguishable from the 1D result, $E_{\mathrm{rel}}^{\mathrm{int}}
\approx 8\times10^{-3}$.  For the point-source case, the plume-core error
should be a few percent at $\Pe_L=20$ and degrade gracefully as $\Pe_L$
increases.  For random $K$, the PINN is expected to be less accurate than the
FV reference; the FNO is the appropriate tool for many-query problems.

\section{Three-dimensional roadmap}
\label{sec:3d}

The 3D plan is conditional and starts with the cheapest model that contains
vertical physics.

\begin{table}[htbp]
\centering
\small
\begin{tabularx}{\textwidth}{@{}lYl@{}}
\toprule
\tabhead{Phase} & \tabhead{Work} & \tabhead{Gate}\\
\midrule
3D-0 & 2.5D axisymmetric $(r,z,t)$ and vertical-slice models; shared 3D
       physics; Thiem/MMS checks. & axisymmetric OK\\
3D-1 & Extruded 3D verification. & G3.1\\
3D-2 & 3D Green's functions, steady and transient. & G3.2\\
3D-3 & 3D Darcy flow, partially penetrating wells; compare to axisymmetric. & G3.3\\
3D-4 & Layered and random 3D $K$; FV reference; MMS; grid convergence. & G3.4\\
3D-5 & Autoregressive FNO3d; rollout; resolution transfer. & G3.5\\
3D-6 & Surrogate-based uncertainty and remediation application. & G3.6\\
\bottomrule
\end{tabularx}
\caption{Three-dimensional work breakdown.  If 3D-0 answers the site question,
stop there and document why full 3D was not needed.}
\label{tab:3d}
\end{table}

\paragraph{The 2.5D bridge.}
The axisymmetric model uses inputs $(r,z,\tau)$ and the operators
\begin{equation}
\label{eq:rz}
\nabla\cdot\vv=\frac1r\pdd{(rv_r)}{r}+\pdd{v_z}{z},
\end{equation}
with a well at $r=r_w>0$.  The residual is weighted by $r$ in the loss to avoid
the coordinate singularity at $r=0$.  The vertical-slice model uses
$(x,z,\tau)$ with uniform $y$ and no-flux conditions at $y=0,L_y$.  Both are
genuine subsets of the 3D equations and much cheaper to train.

\paragraph{When to stop at 2.5D.}
If the question is a single well in a layered aquifer with radial symmetry,
the axisymmetric model is exact and the full 3D model adds cost without adding
information.  Full 3D is justified by lateral heterogeneity, multiple
non-aligned wells, a non-axisymmetric source, or a site geometry that breaks
the symmetry.

\paragraph{3D FNO.}
Start with an autoregressive FNO3d on a $64^3$ grid, 8--12 modes, and 32
channels.  Use a rollout loss and a held-out set of conductivity fields.  Only
after G3.5 passes at $64^3$ should the grid be refined to $128^3$; at that
point use factorized spectral convolutions, mixed precision, gradient
checkpointing, and possibly multiple GPUs.  Storage is the other bottleneck:
one $128^3$ float32 snapshot is about \SI{8.4}{\mega\byte}, so a 100-snapshot
scenario is about \SI{840}{\mega\byte}, and 500 scenarios are about
\SI{420}{\giga\byte}.  Save every $k$-th step, compress, and use Zarr or HDF5.

\section{Computational scaling}
\label{sec:compute}

\begin{table}[htbp]
\centering
\small
\begin{tabularx}{\textwidth}{@{}lYYYY@{}}
\toprule
\tabhead{Model} & \tabhead{Inputs} & \tabhead{Typical points} &
\tabhead{Hardware} & \tabhead{Wall-clock}\\
\midrule
1D PINN & 2 & $1.5\times10^3$ per step & CPU & minutes--hours\\
2D PINN & 3 & $2\times10^4$--$2\times10^5$ & CPU or 8--16 GB GPU & hours--days\\
3D PINN & 4 & $2\times10^5$--$10^6$ & 40--80 GB GPU & days--weeks\\
2D FNO & 2D fields & $128^2$, 500--1000 scenarios & 16--24 GB GPU & 12--48 h\\
3D FNO & 3D fields & $64^3$, 200--500 scenarios & 24--40 GB GPU & 2--7 days\\
\bottomrule
\end{tabularx}
\caption{Order-of-magnitude compute budget.  The 3D PINN and FNO numbers assume
chunked residual evaluation, gradient accumulation, and mixed precision where
possible.}
\label{tab:compute}
\end{table}

The dominant cost in a PINN is the second derivative of the network with
respect to the inputs, evaluated on every collocation point.  The cost scales
linearly with $N_r$ and with the number of network parameters, and the memory
scales with the batch size times the graph depth.  Chunking the residual and
accumulating gradients keeps the memory bounded and is mandatory in 3D.  The
dominant cost in a 3D FNO is the spectral convolution: the FFT is
$\mathcal{O}(N^3\log N)$ per channel, and the weight tensor grows as
$k_xk_yk_zc_{\mathrm{in}}c_{\mathrm{out}}$.  Reducing the number of modes is the
first lever; reducing the number of channels or factorizing the spectral
convolution is the second.

\section{Risks and mitigations}
\label{sec:risks}

\begin{table}[htbp]
\centering
\small
\begin{tabularx}{\textwidth}{@{}YXX@{}}
\toprule
\tabhead{Risk} & \tabhead{Why it happens} & \tabhead{Mitigation}\\
\midrule
Inlet/wall corner singularities & Dirichlet inlet meets a no-flux wall &
report interior metrics; round or exclude the corner; local refinement\\
No-flux hard constraints are difficult & flux conditions are not point values &
zero-normal-derivative mask or soft penalty; monitor $\partial u/\partial n$\\
Coupled flow--transport instability & transport residual contains
$\nabla^2h$ and $\nabla h$ & decouple first; freeze the velocity; try
monolithic later\\
Heterogeneity is high-frequency & FNO truncates modes & more modes; U-NO or
multi-resolution; local PINN refinement\\
High-$\Pe$ fronts are under-resolved & sharp gradients & Fourier features;
causal weighting; time marching; RAR near the front\\
FV reference is not conservative & naive advection/dispersion & TPFA, upwind,
implicit diffusion; mass-balance test every run\\
FNO rollout diverges & autoregressive error growth & pushforward; multi-step
loss; spectral normalization; conservative correction\\
3D memory blow-up & autograd on $10^6$ points & chunking; gradient
accumulation; checkpointing; mixed precision\\
3D FNO parameter explosion & $k^3c_{\mathrm{in}}c_{\mathrm{out}}$ &
8--12 modes; 32 channels; factorized or low-rank spectral convolution\\
Unnecessary 3D & 2.5D would answer the question & written justification before
3D-3; stop at 3D-0 if possible\\
\bottomrule
\end{tabularx}
\caption{Principal risks and the planned mitigations.  The first and last rows
are the ones most likely to waste time.}
\label{tab:risks}
\end{table}

\section{Deliverables and repository layout}
\label{sec:deliverables}

The two plans share a common core.  The expected layout is:
\begin{verbatim}
TESIS/
  physics/
    transport.py       # ADE residual, 2D/3D dispersion tensor
    flow.py            # Darcy residual, wells, K-field generators
    references.py      # OB, Green's functions, MMS (2D)
    references3d.py    # 3D Green's functions, axisymmetric limits, MMS
    fv2d.py            # 2D finite-volume reference solver
    fv3d.py            # 3D finite-volume reference solver
    rz.py              # axisymmetric (r,z) operators
    metrics.py         # rel-L2, mass balance, BTC, plume moments
  solve_2d_pinn.py
  solve_2d_fno.py
  solve_rz_pinn.py
  solve_3d_pinn.py
  solve_3d_fno.py
  configs/2d/*.yaml
  configs/3d/*.yaml
  out/2d/<run_id>/     # config, metrics, model, figures
  out/3d/<run_id>/
  data/3d/*.zarr
  tests/
    2D_plan.md
    3D_plan.md
    2D_results.md
    3D_results.md
\end{verbatim}
The existing \texttt{hydrology.py} and \texttt{solve\_1d\_pinn.py} remain the
trusted regression baseline.  Every run writes a \texttt{run\_config.json}
with the effective configuration, git SHA, device, dtype, and seed, following
the conventions in \texttt{GROWTH.md}.

\section{Summary}
\label{sec:summary}

\begin{table}[htbp]
\centering
\small
\begin{tabularx}{\textwidth}{@{}lYYY@{}}
\toprule
\tabhead{Aspect} & \tabhead{1D (done)} & \tabhead{2D (next)} & \tabhead{3D (conditional)}\\
\midrule
Physics & ADE, uniform flow & ADE + $D_T$, wells, layered $K$ &
3D flow, partial penetration, vertical anisotropy\\
Reference & Ogata--Banks & extruded OB; Green's function; MMS; FV &
extruded OB; 3D Green's function; 2.5D; MMS; FV\\
PINN inputs & $(\xi,\tau)$ & $(\xi,\eta,\tau)$ & $(\xi,\eta,\zeta,\tau)$\\
Hard BC & inlet/outlet/IC & box mask, no-flux walls & 3D box mask, layers\\
FNO & 1D operator & FNO2d, rollout & FNO3d, factorized, rollout\\
Main metric & rel-$L^2$ + interior & + mass balance, BTC, cleanup &
+ vertical profiles, layer BTC\\
Cost & minutes & hours--days & days--weeks\\
\bottomrule
\end{tabularx}
\caption{Summary of the three modelling tiers.  Each tier reuses the previous
tier's verification discipline.}
\label{tab:summary}
\end{table}

The path from one to three dimensions is not a single leap.  It is a sequence
of small, falsifiable steps:
\begin{enumerate}[leftmargin=1.6em,itemsep=2pt]
  \item prove the 2D/3D implementation on an extruded 1D problem;
  \item add transverse dispersion and verify against a Green's function;
  \item add Darcy flow and wells, decoupled first;
  \item add heterogeneity and verify with MMS and a grid-converged FV solver;
  \item only then train the operator-learning surrogate;
  \item enter 3D only after the 2D gates pass, and start with 2.5D.
\end{enumerate}
This is slower than jumping to a full 3D coupled model, but it is the only way
to know \emph{which} ingredient is responsible when the error is large.  The
1D sandbox already demonstrated the payoff: a $1.35\times10^{-2}$ full-grid
error that looked like a bug turned out to be a measure-zero corner
singularity, a diagnosis that was only possible because the reference and the
interior metric were in place.

\bibliographystyle{plainnat}
\begin{thebibliography}{99}

\bibitem[Bear(1972)]{bear1972}
J.~Bear.
\textit{Dynamics of Fluids in Porous Media}.
American Elsevier, New York, 1972.

\bibitem[Jagtap et~al.(2020)]{jagtap2020xpinn}
A.~D. Jagtap, E.~Kharazmi, and G.~E. Karniadakis.
Conservative physics-informed neural networks on discrete domains for
conservation laws: Applications to forward and inverse problems.
\textit{Computer Methods in Applied Mechanics and Engineering},
365:113028, 2020.

\bibitem[Karniadakis et~al.(2021)]{karniadakis2021piml}
G.~E. Karniadakis, I.~G. Kevrekidis, L.~Lu, P.~Perdikaris, S.~Wang, and L.~Yang.
Physics-informed machine learning.
\textit{Nature Reviews Physics}, 3:422--440, 2021.

\bibitem[Kovachki et~al.(2023)]{kovachki2023neural}
N.~Kovachki, Z.~Li, B.~Liu, K.~Azizzadenesheli, K.~Bhattacharya, A.~Stuart,
and A.~Anandkumar.
Neural operator: Learning maps between function spaces with applications to
PDEs.
\textit{Journal of Machine Learning Research}, 24(89):1--97, 2023.

\bibitem[Krishnapriyan et~al.(2021)]{krishnapriyan2021failure}
A.~S. Krishnapriyan, A.~Gholami, S.~Zhe, R.~Kirby, and M.~W. Mahoney.
Characterizing possible failure modes in physics-informed neural networks.
\textit{Advances in Neural Information Processing Systems}, 34, 2021.

\bibitem[LeVeque(2002)]{leveque2002finite}
R.~J. LeVeque.
\textit{Finite Volume Methods for Hyperbolic Problems}.
Cambridge University Press, 2002.

\bibitem[Li et~al.(2021)]{li2021fno}
Z.~Li, N.~Kovachki, K.~Azizzadenesheli, B.~Liu, K.~Bhattacharya, A.~Stuart,
and A.~Anandkumar.
Fourier neural operator for parametric partial differential equations.
\textit{International Conference on Learning Representations (ICLR)}, 2021.

\bibitem[Lu et~al.(2021)]{lu2021deeponet}
L.~Lu, P.~Jin, G.~Pang, Z.~Zhang, and G.~E. Karniadakis.
Learning nonlinear operators via DeepONet based on the universal approximation
theorem of operators.
\textit{Nature Machine Intelligence}, 3:218--229, 2021.

\bibitem[Ogata and Banks(1961)]{ogata1961solution}
A.~Ogata and R.~B. Banks.
\textit{A Solution of the Differential Equation of Longitudinal Dispersion in
Porous Media}.
U.S. Geological Survey Professional Paper 411-A, 1961.

\bibitem[Raissi et~al.(2019)]{raissi2019pinn}
M.~Raissi, P.~Perdikaris, and G.~E. Karniadakis.
Physics-informed neural networks: A deep learning framework for solving forward
and inverse problems involving nonlinear partial differential equations.
\textit{Journal of Computational Physics}, 378:686--707, 2019.

\bibitem[Thiem(1906)]{thiem1906}
G.~Thiem.
\textit{Hydrologische Methoden}.
Gebhardt, Leipzig, 1906.

\bibitem[Wang et~al.(2021)]{wang2021gradient}
S.~Wang, Y.~Teng, and P.~Perdikaris.
Understanding and mitigating gradient flow pathologies in physics-informed
neural networks.
\textit{SIAM Journal on Scientific Computing}, 43(5):A3055--A3081, 2021.

\bibitem[Zheng and Bennett(2002)]{zheng2002applied}
C.~Zheng and G.~D. Bennett.
\textit{Applied Contaminant Transport Modeling}.
Wiley, New York, 2nd edition, 2002.

\end{thebibliography}

\end{document}
